\documentclass{article}

\usepackage[margin=0.85in]{geometry}
\usepackage{amsmath}
\usepackage{booktabs}
\usepackage{cite}
\usepackage{enumitem}
\usepackage{hyperref}

\setlist{nosep, topsep=2pt, leftmargin=1.5em}

\hypersetup{
    colorlinks=true,
    linkcolor=blue,
    citecolor=blue,
    urlcolor=blue
}

\title{\textbf{YZV411E - Big Data Analytics Project Proposal} \\ \large Scalable Flight Delay Propagation Analysis and Adaptive Turnaround Buffer Assessment}
\author{
    \begin{tabular}{cc}
        Aybek Taha At\i lgan & 150200318 \\
        Osmancan Sar\i & 150230728 \\
        Yavuz Burak Yal\c{c}\i n & 150220752
    \end{tabular}
}
\date{October 2026}

\begin{document}

\maketitle

\section*{Abstract}
\textbf{Problem Definition:} Flight disruptions in air transportation rarely stay confined to their point of origin. When an aircraft flies consecutive legs throughout a day, an incoming arrival delay either gets absorbed by the scheduled ground turnaround buffer or propagates into the next leg as a reactionary delay. Typical flight delay studies treat each flight as an independent observation and train basic classifiers for 15-minute delays, ignoring physical aircraft rotations. Conversely, network propagation models in operations research capture delay trees but rely on single-node solvers that collapse under multi-year national flight volumes. 

To bridge this operational and computational gap, this project develops a distributed data pipeline in Apache Spark to evaluate turnaround buffer dynamics across 31.7 million U.S. domestic flights from 2018 through 2022 ($\approx 9.8$ GB raw tabular data). Using distributed chronological windowing over unique aircraft registrations (\texttt{Tail\_Number}), the pipeline links consecutive legs under strict spatial continuity ($Dest_k = Origin_{k+1}$). We evaluate how scheduled ground buffers absorb or amplify delays across major hubs and empirically benchmark distributed PySpark against a single-machine Pandas/DuckDB baseline across scaled data volumes.

\section{Project Aim and Motivation}
Aviation delays generate billions of dollars in annual operating losses. Secondary (reactionary) delays, caused when an aircraft lands late and the scheduled turnaround cannot absorb the deficit, account for a substantial share of total airspace delay minutes. While recent operational literature has begun exploring adaptive ground buffers to curb propagation, evaluating these mechanisms across multi-year nationwide networks remains computationally challenging.

This project delivers two distinct contributions:
\begin{enumerate}
    \item \textbf{Big Data Engineering Track:} We design a distributed PySpark architecture to clean, partition, and window ~32 million flight records. We evaluate the performance trade-offs between uncompressed CSV and columnar Parquet formats, measure the efficiency of hash-repartitioning by \texttt{Tail\_Number}, and conduct systematic benchmarks against single-machine baselines to identify the empirical crossover threshold where distributed scaling justifies cluster overhead.
    \item \textbf{Aviation Operations Track:} We model the empirical relationship between scheduled turnaround duration and downstream delay variation ($\Delta D = \text{DepDelay}_{k+1} - \text{ArrDelay}_k$). We compare uniform buffer schedules against airport-, hour-, and congestion-adaptive alternatives, assessing the trade-off between delay resilience and aircraft fleet utilization (idle ground time).
\end{enumerate}

\noindent\textbf{Why Big Data Technologies Are Required:} Tracking aircraft itineraries across five years requires stateful window operations across 31.7 million rows. A commercial airliner flies 4 to 8 legs daily. Chaining these legs requires a distributed window transformation:
\begin{center}
\texttt{Window.partitionBy("Tail\_Number").orderBy("FlightDate", "CRSDepTime")}
\end{center}
We hypothesize that single-machine in-memory processing will encounter memory exhaustion and heavy swap thrashing during global grouping on 32 million rows. PySpark overcomes this bottleneck through hash-partitioning aircraft across worker cores, spilling shuffle data to disk when heap limits are reached, and leveraging the Catalyst optimizer for column pruning before execution.

\section{Data}
The study uses the multi-year flight dataset curated by Rob Mulla on Kaggle (\url{https://www.kaggle.com/datasets/robikscube/flight-delay-dataset-20182022}), derived from official U.S. Bureau of Transportation Statistics (BTS) Airline On-Time Performance reporting (14 CFR Part 234).
\begin{itemize}
    \item \textbf{Type and Volume:} Five annual files (\texttt{Combined\_Flights\_2018} through \texttt{2022}) totaling \textbf{9.8 GB in raw CSV} ($\approx 31.7$ million records, 62 columns), exceeding the 1 GB course minimum. Snappy-compressed Parquet alternatives occupy $\approx 1.06$ GB total.
    \item \textbf{Core Features:} Airframe identifiers (\texttt{Tail\_Number}, \texttt{Operating\_Airline}), route endpoints (\texttt{Origin}, \texttt{Dest}), local timestamps (\texttt{CRSDepTime}, \texttt{DepTime}, \texttt{CRSArrTime}, \texttt{ArrTime}), and delay metrics (\texttt{DepDelay}, \texttt{ArrDelay}, \texttt{DepDel15}).
    \item \textbf{Preliminary Feasibility Audit (2022 Data):} Our initial audit of 4,078,318 records from 2022 confirms data quality: \texttt{Tail\_Number} missingness is low (0.65\%), while cancellations (3.02\%) and diversions (0.25\%) are explicitly identified. Crucially, consecutive flight sorting shows that \textbf{98.96\%} of same-day pairs satisfy spatial continuity ($Dest_k = Origin_{k+1}$), proving that aircraft rotations can be reliably reconstructed while filtering out irregular ferry flights.
    \item \textbf{Partitioning and Chunking Strategy:} Data is partitioned on disk by \texttt{Year} and \texttt{Month} in columnar Parquet format. In Spark memory, DataFrames are hash-repartitioned by \texttt{Tail\_Number} to isolate complete aircraft trajectories within single executor partitions, minimizing shuffle traffic during lag operations.
\end{itemize}

\section{Tools and Frameworks}
\begin{itemize}
    \item \textbf{Apache Spark (PySpark):} Distributed engine for ingestion, schema casting, hash-repartitioning, and chronological window transformations across multi-core workers.
    \item \textbf{Apache Parquet:} High-efficiency columnar storage with Snappy compression, reducing storage by $\approx 88\%$ and enabling predicate pushdown on \texttt{Origin} and \texttt{FlightDate}.
    \item \textbf{Single-Node Baseline (Pandas / DuckDB):} Used as a control benchmark on 500 MB, 2 GB, and 5 GB slices to quantify speed-up ratios and measure where single-node architectures saturate memory.
    \item \textbf{Plotly / Matplotlib:} Generation of buffer response curves, spatial delay propagation flows, and scalability benchmark profiles.
\end{itemize}

\section{Expected Outcomes and Evaluation}
\begin{enumerate}
    \item \textbf{Distributed Scalability Profiles:} Controlled benchmarks comparing single-node execution against PySpark across data scales (500 MB to 9.8 GB). Metrics include wall-clock runtime, peak heap consumption, and speed-up curves.
    \item \textbf{Turnaround Buffer Response Models:} Empirical characterization of net delay variation $\Delta D = \text{DepDelay}_{k+1} - \text{ArrDelay}_k$ as a function of scheduled buffer $B_{sched} = \text{CRSDepTime}_{k+1} - \text{CRSArrTime}_k$, measuring the transition from delay amplification ($\Delta D > 0$) to delay absorption ($\Delta D \le 0$).
    \item \textbf{Adaptive Buffer Assessment:} Comparing uniform buffer allocations against hub- and time-varying alternatives, quantifying the operational trade-off between delay reduction and aircraft ground idle time.
    \item \textbf{Reproducible Pipeline:} An open-source, documented PySpark repository with modular scripts for ingestion, sequence linking, and performance profiling.
\end{enumerate}

\section{Progress Plan}
Milestones scheduled prior to the progress report deliverable:
\begin{itemize}
    \item \textbf{Weeks 1--2 (Data Ingestion and Pipeline Setup):} Ingest full 2018--2022 files, audit multi-year schema consistency, and construct the partitioned Parquet conversion pipeline.
    \item \textbf{Weeks 3--4 (Exploratory Profiling and Baseline):} Generate baseline delay statistics across carriers, airport hubs, and pre/post-pandemic periods; profile arrival delay distributions.
    \item \textbf{Weeks 5--6 (Sequence Linking Implementation):} Build the PySpark windowing pipeline with operational boundaries: enforce spatial continuity ($Dest_k = Origin_{k+1}$) and classify turnaround times into within-day turns ($30 \le \Delta t \le 240$ min) versus overnight rests ($\Delta t \ge 480$ min).
    \item \textbf{Weeks 7--8 (Benchmarking and Buffer Modeling):} Run standardized scalability benchmarks (Pandas vs. PySpark) across core counts; compute empirical buffer absorption functions.
    \item \textbf{Weeks 9--10 (Synthesis and Final Deliverable):} Evaluate adaptive buffer trade-offs, generate interactive flow visualizations, and author the final project deliverable.
\end{itemize}

\begin{thebibliography}{99}\small
\bibitem{beatty2001}
R.~Beatty et al., ``Preliminary evaluation of flight delay propagation through an airline schedule,'' in \textit{4th USA/Europe ATM R\&D Seminar}, 2001.
\bibitem{pyrgiotis2013}
N.~Pyrgiotis et al., ``A complex network approach to modeling delay propagation in large transportation networks,'' \textit{Transp. Res. Part C}, vol.~27, pp.~208--238, 2013.
\bibitem{fleurquin2013}
P.~Fleurquin et al., ``Systemic delay propagation in the US airport network,'' \textit{Sci. Rep.}, vol.~3, p.~1159, 2013.
\bibitem{bts2023}
Bureau of Transportation Statistics, ``Reporting Carrier On-Time Performance (14 CFR Part 234),'' U.S. Department of Transportation, Tech. Rep., 2023.
\end{thebibliography}

\end{document}